\textbf{1. Input data.} The specified dimensions and sound speed are
\[
\boxed{\begin{aligned}
L_{0}&=10.0\,\mathrm{mm}=0.0100\,\mathrm{m},&
c&=340\,\mathrm{m/s},\\
L_{e}&=420\,\mathrm{mm}=0.420\,\mathrm{m},&
L_{c}&=260\,\mathrm{mm}=0.260\,\mathrm{m}.
\end{aligned}}
\]
The supplied Figure~\ref{fig:test02-whistle} identifies the instrument configuration. The main calculation uses the closed--closed idealization developed in the course notes on pp.~109--113. Those notes allow a small opening for the sound to escape while retaining closed displacement conditions at both ends. The supplied lengths are used directly, so no ruler digitization is required.

\textbf{2. Fully extended whistle length.} The supplied fully extended length is
\[
\boxed{L_{e}=0.420\,\mathrm{m}.}
\]

\textbf{3. Fully contracted whistle length.} The supplied fully contracted length is
\[
\boxed{L_{c}=0.260\,\mathrm{m}.}
\]

\textbf{4. Fully extended travel.} The extension travel is the difference between the two specified overall lengths,
\[
\boxed{dL_{\mathrm{I}}=L_{e}-L_{c}=0.420-0.260=0.1600\,\mathrm{m}.}
\]

\textbf{5. Fully extended cavity length.} Adding the residual cavity to the travel gives
\[
\boxed{L_{\mathrm{I}}=L_{0}+dL_{\mathrm{I}}=0.0100+0.1600=0.1700\,\mathrm{m}.}
\]

\textbf{6. Fully extended frequency.} For the adopted closed--closed idealization, the fundamental frequency from the course notes is
\[
f=\frac{c}{2L}.
\]
Therefore,
\[
f_{\mathrm{I}}=\frac{340}{2(0.170)}=1.000\times10^{3}\,\mathrm{Hz}.
\]
The nearest equal-tempered note is
\[
\boxed{f_{\mathrm{I}}=1.000\,\mathrm{kHz},\qquad \text{nearest note }\mathrm{B5}\ (988\,\mathrm{Hz}).}
\]

\textbf{7. Fifty-percent extension travel.} The travel at the specified half-extension is
\[
\boxed{dL_{\mathrm{II}}=0.50dL_{\mathrm{I}}=0.50(0.160)=0.0800\,\mathrm{m}.}
\]

\textbf{8. Fifty-percent extension cavity length.} Adding the residual cavity gives
\[
\boxed{L_{\mathrm{II}}=L_{0}+dL_{\mathrm{II}}=0.0100+0.0800=0.0900\,\mathrm{m}.}
\]

\textbf{9. Fifty-percent extension frequency.} The closed--closed fundamental is
\[
f_{\mathrm{II}}=\frac{340}{2(0.0900)}=1.889\times10^{3}\,\mathrm{Hz}.
\]
The nearest equal-tempered note is
\[
\boxed{f_{\mathrm{II}}=1.889\,\mathrm{kHz},\qquad \text{nearest note }\mathrm{A\sharp6}\ (1.865\,\mathrm{kHz}).}
\]

\textbf{10. Comment.} Increasing the cavity length lowers the fundamental frequency because the closed--closed relation is inversely proportional to \(L\). The half-extended cavity is shorter than the fully extended cavity, so its frequency is higher, as the two calculated values show. The closed--closed model is the main course-note idealization used here; the Test 2 prompt does not explicitly direct the reader to recall Problem B.3.

\textbf{11. Frequency of D7.} Using equal temperament with the standard \(440\,\mathrm{Hz}\) A4 reference,
\[
\boxed{f_{\mathrm{III}}=440\,2^{29/12}=2.35\times10^{3}\,\mathrm{Hz}.}
\]

\textbf{12. Cavity length for D7.} Applying the closed--closed fundamental relation gives
\[
\boxed{L_{\mathrm{III}}=\frac{c}{2f_{\mathrm{III}}}
=\frac{340}{2(2349.318)}=0.0724\,\mathrm{m}.}
\]

\textbf{13. Extension travel for D7.} The required travel beyond the residual cavity is
\[
\boxed{dL_{\mathrm{III}}=L_{\mathrm{III}}-L_{0}=0.0724-0.0100=0.0624\,\mathrm{m}.}
\]

\textbf{14. Required piston extension.} Relative to the full available travel,
\[
\boxed{100\frac{dL_{\mathrm{III}}}{dL_{\mathrm{I}}}
=100\frac{0.0624}{0.160}=39.0\%.}
\]

\textbf{15. Comment.} The D7 target requires approximately \(39.0\%\) of the available piston travel under the main closed--closed model, so it lies within the specified \(0\)--\(100\%\) range. Figure~\ref{fig:test02-whistle} shows a physical open fipple and a closed piston, which motivates the alternative closed--open model from the course notes on pp.~114--117. Under that model, \(f=c/(4L)\), giving \(5.00\times10^{2}\,\mathrm{Hz}\) at full extension, nearest B4, \(944\,\mathrm{Hz}\) at 50\% extension, nearest \(\mathrm{A\sharp5}\), and a D7 position of \(L_{\mathrm{III}}=36.2\,\mathrm{mm}\), \(dL_{\mathrm{III}}=26.2\,\mathrm{mm}\), or \(16.4\%\) travel. Both boundary models produce a feasible D7 position; the closed--closed values above are the primary course-note idealization.

\clearpage
\matlabcaption{MATLAB calculations for B.5.}
\lstinputlisting[style=hwmatlab]{matlab/B5.m}
